% Original: PDF strana 8, donji slajd.
\slajdid{02-s016}
\begin{frame}[label=02-s016]{Drugi put: komplementarna funkcija}
\setlength{\abovedisplayskip}{3pt}\setlength{\belowdisplayskip}{3pt}
\setlength{\abovedisplayshortskip}{3pt}\setlength{\belowdisplayshortskip}{3pt}
\begin{columns}[c,onlytextwidth]
\begin{column}{.63\textwidth}
% element: 02-s016:e01
\Oznaka{Redovi u kojima je \(F=0\) daju \(\bar F=1\)}\vspace{1mm}
\[\bar F=\bar C\bar B\bar A+\bar C\bar BA+\bar CB\bar A+C\bar B\bar A+CBA\]
% element: 02-s016:e02
\Oznaka{De Morganovo pravilo}\vspace{2mm}
Ako je \(\bar F=P_1+\cdots+P_5\), negiramo ceo zbir:
\[F=\overline{\bar F}=\bar P_1\bar P_2\bar P_3\bar P_4\bar P_5.\]
Svaka negacija proizvoda daje potpuni zbir.
\end{column}
\begin{column}{.30\textwidth}\centering
\Oznaka{\(F\) i komplementna funkcija \(\bar F\)}\vspace{2mm}
{\fontsize{9}{11}\selectfont\begingroup
\setlength{\tabcolsep}{2pt}
\setlength{\aboverulesep}{1pt}\setlength{\belowrulesep}{1pt}
\renewcommand{\arraystretch}{1.0}
\par\noindent
\begin{tabularx}{\linewidth}{*{5}{>{\centering\arraybackslash}X}}
\toprule C&B&A&F&\(\bar F\)\\\midrule
\textcolor{DEcrvena}{0}&\textcolor{DEcrvena}{0}&\textcolor{DEcrvena}{0}&\textcolor{DEcrvena}{0}&\textcolor{DEcrvena}{1}\\
\textcolor{DEcrvena}{0}&\textcolor{DEcrvena}{0}&\textcolor{DEcrvena}{1}&\textcolor{DEcrvena}{0}&\textcolor{DEcrvena}{1}\\
\textcolor{DEcrvena}{0}&\textcolor{DEcrvena}{1}&\textcolor{DEcrvena}{0}&\textcolor{DEcrvena}{0}&\textcolor{DEcrvena}{1}\\
0&1&1&1&0\\
\textcolor{DEcrvena}{1}&\textcolor{DEcrvena}{0}&\textcolor{DEcrvena}{0}&\textcolor{DEcrvena}{0}&\textcolor{DEcrvena}{1}\\
1&0&1&1&0\\
1&1&0&1&0\\
\textcolor{DEcrvena}{1}&\textcolor{DEcrvena}{1}&\textcolor{DEcrvena}{1}&\textcolor{DEcrvena}{0}&\textcolor{DEcrvena}{1}\\\bottomrule
\end{tabularx}\par
\endgroup
}
\end{column}
\end{columns}
% element: 02-s016:e03
\[\textcolor{DEcrvena}{F=(C+B+A)(C+B+\bar A)(C+\bar B+A)(\bar C+B+A)(\bar C+\bar B+\bar A)}\]
\end{frame}
